\subsection{Numerical range}

DEFINITION 2. --- Let E be a complex Hilbert space and let u be an endomorphism of E. We call the numerical range of \(u\) the set of complex numbers of the form \(\langle x|u(x)\rangle\) , where x ranges over the unit sphere of E. We denote it by \(\iota(u)\) .

The numerical range of \(u^{*}\) is the image of \(\iota(u)\) by complex conjugation. For all complex numbers \(\lambda\) and \(\mu\) , the numerical range of \(\lambda u + \mu 1_{\mathrm{E}}\) is equal to \(\lambda \iota(u) + \mu\) .

PROPOSITION 6. — Let E be a complex Hilbert space and let u be an endomorphism of E.

\begin{enumerate}
\def\labelenumi{\alph{enumi})}
\item
  The set of eigenvalues of u is contained in \(\iota(u)\) ;
\item
  The spectrum of u is contained in the closure of \(\iota(u)\) in C.
\end{enumerate}

Let \(\lambda\) be an eigenvalue of \(u\) and let \(x \in \mathrm{E}\) be a nonzero vector such that \(u(x) = \lambda x\) . Replacing \(x\) by \(x / \| x\|\) , we may assume that \(\| x \| = 1\) . Then, \(\langle x | u(x) \rangle = \lambda\) , hence \(\lambda \in \iota(u)\) .

We prove \(b\) ). By considering \(u - \lambda 1_{\mathrm{E}}\) , we reduce to proving that if 0 belongs to the spectrum of \(u\) , then 0 is adherent to \(\iota(u)\) .

Suppose first that there exists a real number \(c > 0\) such that \(\| u(x)\| \geqslant c\) for all \(x\) of norm 1 in E. The endomorphism \(u\) is then injective and has closed range (Lemma 8 of I, p.~107). Since it is not invertible by hypothesis, it is not surjective. Consequently, the orthogonal complement of the kernel of \(u^{*}\) is not equal to E (EVT, V, p.~41, prop. 4), which proves that the kernel of \(u^{*}\) is nonzero. Thus 0 belongs to \(\iota(u^{*})\) , hence to \(\iota(u)\) .

If the preceding hypothesis is not valid, then for every integer \(n \geqslant 1\) , there exists a vector \(x_{n}\) of norm 1 in E such that \(\|u(x_{n})\| \leqslant 1/n\) . We then have \(|\langle x_{n}|u(x_{n})\rangle| \leqslant 1/n\) , which implies that 0 belongs to the closure of \(\iota(u)\) .

PROPOSITION 7 (Hausdorff-Toeplitz Theorem)

Let E be a complex Hilbert space and let \(u \in \mathcal{L}(\mathrm{E})\) . The numerical range \(\iota(u)\) is a convex subset of C.

We will need two lemmas to prove this proposition.

Lemma 5. --- Let E be a complex Hilbert space of dimension 2. Equip the real vector space \(\mathcal{L}(\mathrm{E})_h\) of Hermitian endomorphisms of E with the pre-Hilbert norm \(u \mapsto \mathrm{Tr}(u^* u)^{1/2}\) . The set of rank-1 orthogonal projectors of E is the sphere S of radius \(1/\sqrt{2}\) centered at \(\frac{1}{2} 1_{\mathrm{E}}\) in the 3-dimensional affine subspace of endomorphisms of trace 1 in \(\mathcal{L}(\mathrm{E})_h\) .

Let F be the real affine subspace of \(\mathcal{L}(\mathrm{E})_h\) consisting of elements of trace 1. The rank-1 orthogonal projectors of E belong to F (lemma 3, (ii)).

Let \(u \in \mathrm{F}\) . We have \(\| u - \frac{1}{2} 1_{\mathrm{E}} \|^{2} = \mathrm{Tr}(u^{2} - u + \frac{1}{4}) = \mathrm{Tr}(u^{2}) - \frac{1}{2}\) . Consequently, \(u \in \mathrm{S}\) if and only if \(\mathrm{Tr}(u^{2}) = 1\) . Since \(2\det(u) = \mathrm{Tr}(u)^{2} - \mathrm{Tr}(u^{2}) = 1 - \mathrm{Tr}(u^{2})\) , this condition is equivalent to \(\det(u) = 0\) . By the Hamilton--Cayley theorem (A, III, p.~107, prop. 20), we therefore have \(u \in \mathrm{S}\) if and only if \(u^{2} - u = 0\) , which means that \(u\) is a rank-1 Hermitian projector (loc. cit.), whence the result.

Lemma 6. --- Let E be a real normed vector space, let F be a real vector space, and let u: E → F be a non-injective affine map. Let B be a ball in E and S the corresponding sphere. We have \(u(S) = u(B)\) , and in particular, \(u(S)\) is convex.

We reduce to the case where \(u\) is linear and where B is the unit ball of E. We have \(u(S) \subset u(B)\) . Conversely, let \(x \in B\) and let \(y\) be a nonzero element of \(\text{Ker}(u)\) . The image of the continuous map \(t \mapsto \|x + ty\|\) of \(\mathbf{R}\) into \(\mathbf{R}_+\) is an unbounded interval containing the real number \(\|x\| \leqslant 1\) . Hence there exists \(t \in \mathbf{R}\) such that \(\|x + ty\| = 1\) . We then have \(x + ty \in S\) and \(u(x + ty) = u(x)\) , hence \(u(x) \in u(S)\) .

Let us prove Prop. 7. Let x and y be elements of the unit sphere of E; we prove that the segment with endpoints \(\langle x|u(x)\rangle\) and \(\langle y|u(y)\rangle\) is contained in the numerical range of u.

Let F be the subspace of E spanned by x and y. If \(\dim(\mathrm{F}) = 1\) , one has \(\langle x|u(x)\rangle = \langle y|u(y)\rangle\) , whence the assertion. Otherwise, we have \(\dim(\mathrm{F}) = 2\) ; let p be the orthogonal projector of E with image F, and denote by \(u_{F}\) the endomorphism of F given by \(x \mapsto p(u(x))\) . Since p is Hermitian (Lemma 3 of I, p.~133), we have \(\langle z|u_{\mathrm{F}}(z)\rangle = \langle z|u(z)\rangle\) for all \(z \in F\) , so that \(\iota(u_{\mathrm{F}}) \subset \iota(u)\) . We may therefore assume that E = F.

For every element \(z\) of E, let \(v_{z}\) be the Hermitian endomorphism of E defined by \(t \mapsto \langle z|t\rangle z\) ; we have \(\langle z|u(z)\rangle = \mathrm{Tr}(u \circ v_{z})\) . When \(z\) traverses the unit sphere of E, \(v_{z}\) describes the set of rank-1 orthogonal projectors of E, which is a sphere S in the real affine subspace V of \(\mathcal{L}(\mathrm{E})\) formed by the Hermitian endomorphisms of E with trace 1 (Lemma 5). The numerical range of E is therefore the set of \(\mathrm{Tr}(u \circ v)\) , for \(v \in \mathrm{S}\) . The map \(v \mapsto \mathrm{Tr}(u \circ v)\) of V in \(\mathbf{C}\) is linear. Since \(\dim_{\mathbf{R}}(\mathrm{V}) = 3 > \dim_{\mathbf{R}}(\mathbf{C})\) , it is not injective; it then follows from Lemma 6 that \(\iota(u)\) is convex.

